\scanpage{26}
\mspage{24}
% archive sheet number: 25
% hand: Grothendieck
% sheet_kind: cancelled_draft
% element: text | cancelled upper block, lines 1–7 | 5.5 On en conclut
\deleted{5.5. On en conclut que \deleted{\unsure{il existe le groupe}}\inserted{l'ens \inserted{$U$}} des $g\in\pi$ tels que $u(g)$ soit unipotent est un sous-groupe ouvert invariant de $\pi$, à savoir $u^{-1}(G^\circ(R))$; c'est évidemment le plus grand sous-groupe \deleted{ouvert} (ou le plus grand sous-groupe ouvert) de $\pi$ \deleted{qui \unsure{y}} opère de façon unipotente. Rappelons que}
% element: display | cancelled upper block, line 8 | pi/U isom (G/G^0)(R)
\[
\deleted{\pi/U\simeq(G/G^\circ)(R)\simeq G(R)/G^\circ(R).}
\tag*{(5.5.1)}
\]
% element: text | cancelled middle, lines 9–11 | Soit U' un sous-groupe
\deleted{Soit $U'$ un sous-groupe \deleted{\unsure{d'indice|fini}}\inserted{(ouvert)} de $\pi$ \deleted{qui opère} de façon unipotente, i.e. un ss-groupe ouvert de $U$, je dis qu'on a alors}
% element: display | cancelled middle, line 12 | E^U' = E^U = E^G^0
\[
\deleted{E^{U'}=E^U=E^{G^\circ}.}
\]
% element: text | cancelled middle-lower, lines 13–17 | En effet
\deleted{En effet, on peut supposer que $U'$ est invariant; alors $U/U'$ opère sur $E^{U'}$, et \deleted{\unsure{et qui}}\inserted{c'est} un groupe fini qui opère de façon unipotente, donc trivialement, donc $E^{U'}=E^U$. On a $E^U=E^{G^\circ}$, car on a déjà vu que $\overline{u(U)}=G^\circ$.}

% element: text | lower, line 18 and interline | 5.6 Supposons maintenant
5.6. Supposons maintenant que $R=\mathbb Q_\ell$, \deleted{Alors}\inserted{et $G^\circ$ comm.} \inserted{i.e. que l'opération de $\pi$ sur $E$ soit ``quasi-abélienne''.}
% element: margin note | left lower margin, beside line 19 | Donc G^0=G_a^d
% slanted
\margin{Donc $G^\circ\simeq\unsure{\mathbb G_a^d|G_a^d}$. (avec structure vectorielle canonique}
% element: text | lower, lines 19–21 | Alors U to G^0(Q_l)
Alors $U\longrightarrow G^\circ(\mathbb Q_\ell)\simeq\mathbb Q_\ell^d$ a comme image un sous-groupe compact de $\mathbb Q_\ell^d$, lequel est nécessairement isomorphe à un $\mathbb Z_\ell^{d'}$ $(d'\leq d)$.

% element: text | bottom, lines 22–26 | Comme u(U) engendre
Comme \deleted{l'image de $U$} $u(U)$ engendre le le groupe algébrique $G^\circ$, il s'ensuit aussitôt que $d'=d$, de façon précise $u(U)$ est un sous-$\mathbb Z_\ell$-module de $G^\circ(\mathbb Q_\ell)$ qui engendre l'espace vectoriel $G^\circ(\mathbb Q_\ell)$.
